\documentclass[sigconf,nonacm]{acmart}
\settopmatter{printacmref=false}
\renewcommand\footnotetextcopyrightpermission[1]{}
\pagestyle{plain}

\begin{document}

\title{Expense Tracker: A Personal Spending Manager Built with Jetpack Compose and XML Views}

\author{Saad}
\affiliation{\institution{Faculty of Computers and Information Sciences, Ain Shams University}\city{Cairo}\country{Egypt}}
\email{saad@example.com}

\author{Youssef}
\affiliation{\institution{Faculty of Computers and Information Sciences, Ain Shams University}\city{Cairo}\country{Egypt}}
\email{youssef@example.com}

\author{Ashraf}
\affiliation{\institution{Faculty of Computers and Information Sciences, Ain Shams University}\city{Cairo}\country{Egypt}}
\email{ashraf@example.com}

\author{Moaaz}
\affiliation{\institution{Faculty of Computers and Information Sciences, Ain Shams University}\city{Cairo}\country{Egypt}}
\email{moaaz@example.com}

\author{Gamil}
\affiliation{\institution{Faculty of Computers and Information Sciences, Ain Shams University}\city{Cairo}\country{Egypt}}
\email{gamil@example.com}

\author{Shalaby}
\affiliation{\institution{Faculty of Computers and Information Sciences, Ain Shams University}\city{Cairo}\country{Egypt}}
\email{shalaby@example.com}

\begin{abstract}
Many students and young professionals lose track of where their money goes because they
record spending late, or not at all. This paper presents \emph{Expense Tracker}, a lightweight
native Android application that lets users record, categorize, edit, share, and delete daily
expenses and compare their total spending against a monthly budget. The application is
written in Kotlin and combines two Android UI toolkits: the main expense list and the entry
dialog are built with Jetpack Compose, while the details screen and the budget banner are built
with traditional XML layouts. Data is stored locally with the Room library so the app works
offline, and navigation between screens relies on explicit and implicit Intents. We describe the
problem, the architecture, the implementation, and the lessons learned.
\end{abstract}

\keywords{Android, Jetpack Compose, XML layouts, Room, Kotlin, Personal finance}

\maketitle

\section{Introduction}
Personal finance habits start with awareness. Small daily purchases such as food, transport, and
study materials add up quickly, yet most budgeting applications are complex, require accounts,
or need a constant Internet connection. Our goal was a small, fast, offline-first application
focused on one job: making it effortless to log an expense and immediately see how it affects the
monthly budget. The app was developed for the Mobile Computing course (FNI-421) to practice
modern Android development.

\section{Problem Statement and Objectives}
Users need a quick way to record expenses, classify them, and see their total spending. The
objectives are: (1) full create, read, update, and delete (CRUD) operations on expenses;
(2) persistent on-device storage; (3) a clean and responsive interface; (4) interoperability
between Jetpack Compose and XML views; and (5) clean Kotlin code that favors immutable values,
lambdas, and data classes.

\section{System Design}
\subsection{Architecture}
The app follows a modular MVVM structure separated into \texttt{data/} and \texttt{ui/} packages.
The \texttt{Expense} entity holds a title, amount, category, optional note, and timestamp, while
\texttt{Categories} defines the supported categories and monthly budget. \texttt{ExpenseDao}
exposes suspend functions for insert, update, and delete, and a Kotlin \texttt{Flow} that emits
the list sorted by date whenever the table changes. \texttt{ExpenseRepository} abstracts database
access, and \texttt{ExpenseViewModel} converts this flow into a \texttt{StateFlow} observed by
the interface in a lifecycle-aware way, so the screen always reflects the database.

\subsection{User Interface}
The home screen (\texttt{ExpenseScreen}) uses Jetpack Compose: a Material 3 scaffold, a scrolling
list of \texttt{ExpenseCard} items, a floating action button, and \texttt{AddExpenseDialog} with
input validation and a reusable \texttt{CategoryChipRow}. At the top, a budget banner defined in
XML (\texttt{view\_banner.xml}) is embedded into the Compose hierarchy through \texttt{AndroidView}
inside \texttt{SpendingBanner}; it shows the total spent and a progress bar against the monthly
budget. Selecting an expense opens \texttt{ExpenseDetailActivity}, whose interface is entirely
XML (\texttt{activity\_expense\_detail.xml}) and uses Material text fields, a category spinner,
and buttons.

\subsection{Navigation with Intents}
The list opens the details screen with an explicit Intent that carries the expense identifier.
The ``Share expense'' button sends a plain-text summary through an implicit \texttt{ACTION\_SEND}
Intent, so it can be shared with any installed messaging application.

\section{Persistence}
Room provides compile-time-verified SQL and an abstraction over SQLite. The database is a
singleton, and all queries run off the main thread using coroutines to keep the interface
responsive. Both UI toolkits operate on the same data, so a change made in the XML screen is
reflected immediately in the Compose list.

\section{Implementation Notes}
The code is split into small single-responsibility files across \texttt{data/},
\texttt{ui/screens/}, \texttt{ui/components/}, and \texttt{ui/detail/}. Immutable \texttt{val}
declarations are preferred, \texttt{var} is used only for local UI state, and lambdas are passed
down the composable tree (state hoisting). Amounts are validated as positive decimal numbers
before they are saved, and an empty-state message guides first-time users.

\section{Challenges and Lessons Learned}
Mixing Compose and XML required a shared Material 3 theme so colors and typography stay
consistent. Embedding a View in Compose clarified the difference between the \texttt{factory}
block, which inflates once, and the \texttt{update} block, which reruns on every state change.
Dividing work by layer across six team members (database, MVVM/navigation, Compose list, Compose
form dialog, XML-in-Compose banner + docs, and full XML details screen) allowed everyone to
contribute through separate Git branches and pull requests.

\section{Conclusion and Future Work}
Expense Tracker shows that Compose and XML can coexist in a single coherent application and that
Room makes local persistence simple. Future work includes charts per category, date filters,
exporting to CSV, and optional cloud backup.

\begin{acks}
We thank Dr.\ Soha Nabil Hussein for guidance throughout the Mobile Computing course.
\end{acks}

\end{document}
